% 19 November 2026. Drafted for Edward Baker with OpenAI Codex assistance.
% E17 helper explicitly dispatched gpt-6.1-sol/ultra; backend unattested.
\section{A common dyadic phase still leaves an envelope obstruction}
\label{sec:dyadic-group-obstruction}

We strengthen Section~\ref{sec:arithmetic-obstruction} by grouping all
frequencies with the same odd core and retaining their common dyadic
phase across groups. The remaining full-circle envelope still has a
negative complete allowance. All source quantities $q=f-c$, $v$,
$D_\kappa$, $E$, $\beta_m,a_m$, $h_m$ and the segment $\mathcal F_0$
are exactly those of \eqref{af:definitions}--\eqref{af:segment}.
In particular, both full term-error sums and both retained delta payments
remain in $v$, and $|H/\mathcal B-q|\le v$. No new source estimate is
imported.

\subsection{Exact dyadic groups and their allowance}

Write $\rho=\gamma E/\bar E=g e^{\ii\theta}$ as before. Let
$\mathcal R_N$ be the positive odd integers at most $30N$, and put
\[
 I_1=\{j\ge1:2^j\le30N\},\qquad
 I_r=\{j\ge0:r2^j\le30N\}\quad(r\ge3\text{ odd}).
\]
For real auxiliary phases define the finite polynomials and groups
\begin{align}
 B_r(u)&=\sum_{j\in I_r}\beta_{r2^j}2^{-j\sigma}e^{\ii ju},&
 A_r(u)&=\sum_{j\in I_r}a_{r2^j}2^{-j\sigma}e^{\ii ju},\notag\\
 G_r(\eta,u)&=r^{-\sigma}
    \bigl(e^{\ii\eta}B_r(u)+g e^{-\ii\eta}\overline{A_r(u)}\bigr).
 \label{dg:groups}
\end{align}
With $F=S+\gamma\bar T$, the exact finite convolution gives
\begin{equation}
 e^{-\ii\theta/2}EF=Z_0+
 \sum_{r\in\mathcal R_N}G_r(\eta_r,u),\qquad
 Z_0=e^{-\ii\theta/2}+g e^{\ii\theta/2},\quad
 u=-\tau\log2,\quad\eta_r=-\tau\log r-\theta/2.
 \label{dg:identity}
\end{equation}
Every $m\ge2$ has its unique dyadic representation; the excluded
$r=1,j=0$ term is the displayed anchor. All cutoff terms through $30N$
are retained, without assuming multiplicativity of $b_n$.

Use $\|w\|_g^2=(\Re w/(1+g))^2+(\Im w/(1-g))^2$.
Its anchor norm is one and $|w|\ge(1-g)\|w\|_g$. Set
\begin{align}
 R_r(u)&=(1-g)\sup_{\eta\in[0,2\pi]}\|G_r(\eta,u)\|_g,\notag\\
 C_{\rm dyad}&=\sup_{u\in[0,2\pi]}\sum_rR_r(u),&
 C_{\rm sep}&=\sum_r\sup_{u\in[0,2\pi]}R_r(u),\notag\\
 K_{\rm dyad}&=1-g-C_{\rm dyad},&
 \ell_{\rm dyad}&=K_{\rm dyad}/|E|-D_\kappa-|c|,\qquad
 m_{\rm dyad}=\ell_{\rm dyad}-v.
 \label{dg:allowance}
\end{align}
The coefficients, $N,\sigma,g$ are frozen at the actual observation
while these auxiliary suprema are taken. Only the core rotations are
independent; the same $u$ is retained across all groups.

For a real matrix, $\|\cdot\|_2$ denotes its Euclidean operator norm.
Direct expansion of $G_r$ in $\cos\eta,\sin\eta$ gives
\begin{equation}
 M_r(u)=\begin{pmatrix}
 \dfrac{\Re B_r+g\Re A_r}{1+g}&
       -\dfrac{\Im B_r+g\Im A_r}{1+g}\\[4pt]
 \dfrac{\Im B_r-g\Im A_r}{1-g}&
        \dfrac{\Re B_r-g\Re A_r}{1-g}
 \end{pmatrix},\qquad
 R_r(u)=(1-g)r^{-\sigma}\|M_r(u)\|_2.
 \label{dg:matrix}
\end{equation}
Indeed the scaled real coordinates are
$r^{-\sigma}M_r(u)(\cos\eta,\sin\eta)^{\mathsf T}$, whose supremum
is exactly that operator norm. Reverse triangle in \eqref{dg:identity}
therefore gives $|EF|\ge1-g-\sum_rR_r(u)\ge K_{\rm dyad}$.
Paying \eqref{af:kappa} and $|c|$ proves $\ell_{\rm dyad}\le|q|$,
including when $K_{\rm dyad}<0$. The divisor remains the exact $|E|$.

Each $j$th matrix coefficient in \eqref{dg:matrix} is a diagonal
matrix followed by rotation through $ju$, with norm
$2^{-j\sigma}h_{r2^j}/(1-g)$. Its triangle inequality and the unique
dyadic partition imply
\begin{equation}
 C_{\rm dyad}\le C_{\rm sep}\le
 C_{\rm ind}:=\sum_{m=2}^{30N}h_m m^{-\sigma},\qquad
 m_E\le m_{\rm dyad}\le m_{\rm dir}=|f-c|-v.
 \label{dg:comparison}
\end{equation}
Direct evaluation receives every polynomial, actual phase and complete
payment and may use this same grouping. No cost advantage is asserted.

\subsection{A Fourier mean forces the same negative constant}

\begin{theorem}[Scoped common-phase envelope obstruction]
On $\mathcal F_0$ in \eqref{af:segment}, including both one-sided cutoff
references, $\ell_{\rm dyad}<-1/400$ and $m_{\rm dyad}<-1/400$.
The scope is the fixed three-prime multiplier, anchor one, odd-core
dyadic groups, independent core rotations, one common full-circle
dyadic-phase envelope, and separate kappa, correction and source costs.
\end{theorem}

\begin{proof}
Write $\langle\cdot\rangle=(2\pi)^{-1}\int_0^{2\pi}(\cdot)\,du$.
For odd $r\ge3$, the real constant coefficients of the finite
polynomials give
\[
 \langle M_r\rangle=\operatorname{diag}
 \left(\frac{\beta_r+ga_r}{1+g},
       \frac{\beta_r-ga_r}{1-g}\right).
\]
The integral triangle inequality, applied to $M_r$ on each unit vector,
proves $\|\langle M_r\rangle\|_2\le\langle\|M_r\|_2\rangle$.
Thus \eqref{dg:matrix} implies
\[
 \langle R_r\rangle\ge h_r r^{-\sigma},\qquad
 C_{\rm dyad}\ge\left\langle\sum_rR_r\right\rangle
 \ge\sum_{\substack{r\ge3\ {\rm odd}\\r\le30N}}h_r r^{-\sigma}.
\]
The core $r=1$ has zero mean matrix and nonnegative cost. No limiting
argument or nonvanishing assumption enters these finite norm identities.

Section~\ref{sec:arithmetic-obstruction} proves uniformly on
$\mathcal F_0$ that $N\ge33$, $3/5<\sigma<4/5$, $0<g<3/5$,
$0<|E|<8$, and the eight-prime mass \eqref{af:prime-mass} exceeds
$21/20$. Each prime there is an odd core with
$\beta_p=b_p$, $a_p=p^{1/5}b_p$, hence $h_p\ge(1-g)b_p$.
Consequently
\[
 C_{\rm dyad}>\frac{21}{20}(1-g),\qquad
 K_{\rm dyad}<-\frac{1-g}{20}<-\frac1{50},\qquad
 \frac{K_{\rm dyad}}{|E|}<-\frac1{50|E|}<-\frac1{400}.
\]
The three remaining payments are nonnegative. The bound eight is used
only to upper-bound this already-defined negative allowance.
\end{proof}

For this estimator a positive complete allowance would require exactly
\[
 C_{\rm dyad}<1-g-|E|(D_\kappa+|c|+v)\le1-g,
\]
which the mean bound prevents. Clipping $\ell_{\rm dyad}$ to zero
cannot beat $v\ge0$ strictly. These statements upper-bound a lower
allowance, not $|q|$ or $|H/\mathcal B|$.
The smaller actual-phase cost $\sum_rR_r(u(x))$ also gives a valid
allowance, whose sign is undetermined here. A circle mean is not a
pointwise lower bound at $u(x)$, and no adverse maximizing phases are
claimed to occur at one physical height.

\subsection{Interfaces and the information still missing}

The bounds $25/8<\pi<22/7$ and $x_0<14001$ give
$x_{33}<14000<x_0<x_{34}$ and $x_{48}<30000<x_{49}$, where
$x_j=4\pi j^2-\pi/40$. Thus the natural indices are $33,\ldots,48$,
with internal seams $j=34,\ldots,48$. Apply each finite formula on
its half-open natural cell and old one-sided limit; the new index owns
the seam and $30000$ is included with $N=48$. Both references retain
all eight primes. A fixed-reference transfer still pays the exact
signed jump and full error in \eqref{uc:seams}; no floor derivative
or unpaid gluing is used.

Failure at the included $x_0$ cannot be repaired by moving a tail join.
The rest of the two-dimensional finite strip, effective tail onset and
overlap, initial verification and full intermediate barrier remain
unacquired. The collision screen in Section~\ref{sec:arithmetic-obstruction}
also applies: selected arithmetic or dyadic identities do not supply the
complete literal theta source or a strict margin. Under any honest
enclosure an included zero forces $\ell-v\le0$; the displayed quadratic
control has perfect approximation but lacks theta arithmetic.

This exclusion leaves actual-height joint phase restrictions, correlations
among cores or with the anchor and correction, other groupings and
products, and contour or heat arguments open. Such an escape needs an
independently checked bound crossing the full paid threshold, with
coefficient variation and every cutoff included and an equally informed
direct comparator. No positive strip floor or Newman bound is acquired.
